% TOP_PROBLEM: {"id": 2100513, "title": "Open Problems in Integrable Systems — Nijenhuis operators: singular points and global properties", "category_id": 11, "category": {"id": 11, "name": "dynamical_systems", "display_name": "Dynamical Systems", "description": "Problems about long-term behavior of deterministic systems, Hamiltonian dynamics, and periodic orbits.", "slug": "dynamical-systems", "order_index": 11, "created_at": "2026-07-31T15:26:25.671Z"}, "status": "open"}
% FIRST_POSTED: null
% ENTRY_KIND: statement_only
% Imported from: batch09.tex; UnsolvedMath ID 2100513
\documentclass[11pt]{article}
\usepackage[margin=25mm]{geometry}
\usepackage{fontspec}
\setmainfont{Latin Modern Roman}
\usepackage{amsmath,amssymb,mathtools}
\usepackage{unicode-math}
\setmathfont{Latin Modern Math}
\usepackage{xurl,hyperref}
\hypersetup{colorlinks=true,urlcolor=blue}
\setcounter{secnumdepth}{0}
\setlength{\parindent}{0pt}
\setlength{\parskip}{5pt}
\setlength{\emergencystretch}{3em}
\newcommand{\sref}[2]{\hyperlink{source:#1:#2}{[S#2]}}
\newcommand{\cataloguescope}[1]{\paragraph{Recorded target.}#1}
\begin{document}
% UnsolvedMath ID 2100513; source catalogue code AMR-020-0513
\setcounter{section}{2100512}
\section{Open Problems in Integrable Systems — Nijenhuis operators: singular points and global properties}\label{2100513}
{\small\sffamily UnsolvedMath ID 2100513 · Dynamical Systems}
\subsection{Definitions and mathematical statement}

\paragraph{Shared notation.}$\mathbb N=\{1,2,\ldots\}$, and $\mathbb Z,\mathbb Q,\mathbb R,\mathbb C$ denote the integers, rationals, reals and complex numbers. Any other conventions are specified locally.


For $\sigma=(\sigma_1,\ldots,\sigma_n)$, define the companion matrix $C(\sigma)$ by
\[ C_{i1}=-\sigma_i\quad(1\leq i\leq n),\qquad C_{i,i+1}=1\quad(1\leq i<n), \]
with all remaining entries zero. For a map $\Phi=(\sigma_1,\ldots,\sigma_n)$ in coordinates $x_1,\ldots,x_n$, its Jacobian is $J_{ij}=\partial\sigma_i/\partial x_j$.\par Let $\sigma_i\in\mathbb R[x_1,\ldots,x_n]$ be homogeneous of degree $i$ for $1\leq i\leq n$, and assume they are algebraically independent. Then $D=\det J$ is a nonzero polynomial and
\[ L=\frac{\operatorname{adj}(J)C(\sigma)J}{D}. \]
Each entry is a homogeneous rational function of degree one. Write $Q_{ij}$ for the corresponding numerator in $\operatorname{adj}(J)C(\sigma)J$.
\paragraph{Statement recorded in the supplied source.}
Classify all such collections $(\sigma_1,\ldots,\sigma_n)$ for which $D$ divides every $Q_{ij}$ in $\mathbb R[x_1,\ldots,x_n]$. Equivalently, require every entry of $L$ to be a homogeneous linear polynomial in $x_1,\ldots,x_n$. The reconstructed field then has zero Nijenhuis torsion and corresponds to a left-symmetric algebra, as described in the source. \sref{2100513}{2}
\subsection{Short English statement}
Classify algebraically independent homogeneous coefficient polynomials for which the reconstructed tensor has linear polynomial entries.
\subsection{Sources}
\begin{description}
\item[\hypertarget{source:2100513:1}{S1}] ULAM.AI and the UnsolvedMath Contributors, UnsolvedMath: A Curated Collection of Open Mathematics Problems, record 2100513 (AMR-020-0513). CC BY 4.0. \url{https://huggingface.co/datasets/ulamai/UnsolvedMath}.
\item[\hypertarget{source:2100513:2}{S2}] Alexey Bolsinov, Vladimir S. Matveev, Eva Miranda and Serge Tabachnikov, Open Problems, Questions, and Challenges in Finite-Dimensional Integrable Systems (2018), Problem 5.16 and its complete preceding definitions. Original question and full discussion. \url{https://arxiv.org/abs/1804.03737}.

\end{description}
\label{end:2100513}
\end{document}
